\documentclass[11pt]{article}
\usepackage[margin=1in]{geometry}
\usepackage{amsmath}
\usepackage{graphicx}
\usepackage{hyperref}

\title{Paper Title}
\author{Author Name \\ Affiliation \\ \texttt{email@example.com}}
\date{\today}

\begin{document}
\maketitle

\begin{abstract}
State the problem, the approach and the main result in three or four sentences.
\end{abstract}

\section{Introduction}
\label{sec:intro}
Explain why the problem matters and what is missing in existing work.
Summarize the contributions of this paper.

\section{Method}
\label{sec:method}
Describe the approach. Equations are numbered automatically:
\begin{equation}
  \label{eq:objective}
  \mathcal{L}(\theta) = \frac{1}{N} \sum_{i=1}^{N} \ell\left(f_\theta(x_i), y_i\right).
\end{equation}
Refer to it as Equation~\eqref{eq:objective}.

\section{Results}
\label{sec:results}
Report the main findings and compare them with the baselines.

\section{Conclusion}
\label{sec:conclusion}
Summarize what was shown and what remains for future work.

\end{document}
